\fig{em2:charge-nodes}{电容重新连接前后：虚线框标出应守恒总电荷的上导体}{
\begin{tikzpicture}[scale=.82]
\begin{scope}
\node at(1.5,3){连接前};
\draw[world](0,2.3)--(1,2.3) (2,2.3)--(3,2.3);
\draw[world](0,2.3)--(0,1.35) (0,1.05)--(0,0)
 (3,2.3)--(3,1.35) (3,1.05)--(3,0) (0,0)--(3,0);
\draw[world](-.35,1.35)--(.35,1.35) (-.35,1.05)--(.35,1.05)
 (2.65,1.35)--(3.35,1.35) (2.65,1.05)--(3.35,1.05);
\draw[world](1,2.3)--(1.8,2.65);
\node at(1.5,1.85){$S$};
\node[left]at(-.35,1.2){$C_1$};
\node[right]at(3.35,1.2){$C_2$};
\node[above]at(.1,2.3){$q_{10}$};
\node[above]at(2.9,2.3){$q_{20}$};
\end{scope}
\draw[flow](4,1.3)--(4.8,1.3);
\begin{scope}[xshift=6cm]
\node at(1.5,3.45){闭合后与外界绝缘};
\draw[blue,dashed,rounded corners](-.5,1.2)rectangle(3.5,2.7);
\draw[world](0,2.3)--(3,2.3)
 (0,2.3)--(0,1.35) (0,1.05)--(0,0)
 (3,2.3)--(3,1.35) (3,1.05)--(3,0) (0,0)--(3,0);
\draw[world](-.35,1.35)--(.35,1.35) (-.35,1.05)--(.35,1.05)
 (2.65,1.35)--(3.35,1.35) (2.65,1.05)--(3.35,1.05);
\node at(1.5,2.95){$N:\ Q=q_{10}+q_{20}$};
\node[left]at(-.35,1.2){$C_1$};
\node[right]at(3.35,1.2){$C_2$};
\node[below]at(1.5,0){$M$};
\end{scope}
\end{tikzpicture}}

图\ref{em2:charge-nodes}中，闭合$S$以后两个上极板与连线组成同一个孤立导体$N$。电荷可以在两块上极板之间转移，但不能越过电容绝缘层离开$N$。因此应守恒的是框内总电荷，而非各块极板原有电荷。以下用$U=V_N-V_M$计算共同电压，初始电荷均按各上极板的有向电荷代入；闭合过程保留微小电阻以保证最终稳定。
